% Source: https://artofproblemsolving.com/community/c3h2018256_area_of_triangle
% Stable identity: aops-topic-2018256-q01
Let the square have side $s$. Since $CE=CF$, $CD=CB$ and $\angle D=\angle B=90^\circ$, the right triangles $CDE$ and $CBF$ are congruent (hypotenuse--leg), so $DE=BF=t$. Then $AE=AF=s-t$ and $GC=t$, so $DG=s-t=AF$: the point $G$ lies directly below $F$, and $FG=s$ is perpendicular to $DC$.

\begin{asy}
usepackage("amsmath");
size(5.5cm);
defaultpen(fontsize(10pt));
real s=6, t=s*(2-sqrt(3));
pair D=(0,0), C=(s,0), B=(s,s), A=(0,s);
pair E=(0,t), F=(s-t,s), G=(s-t,0);
fill(D--F--G--cycle,palegreen); fill(A--E--F--cycle,paleyellow);
draw(A--B--C--D--cycle,linewidth(0.9)); draw(E--F--C--cycle,blue); draw(D--F); draw(F--G,dashed);
dot(E);dot(F);dot(G);
label("$A$",A,NW); label("$B$",B,NE); label("$C$",C,SE); label("$D$",D,SW);
label("$E$",E,W); label("$F$",F,N); label("$G$",G,S);
label("$t$",(0,t/2),W); label("$t$",(s-t/2,s),N); label("$t$",(s-t/2,0),S);
label("$s-t$",(0,(s+t)/2),W); label("$s-t$",((s-t)/2,s),N); label("$s$",(s+0.2,s/2),E);
label("\mbox{}",(min(currentpicture,user=true)+max(currentpicture,user=true))/2,black); // keeps the text after the figure black when LaTeX typesets the labels inline
\end{asy}

\textbf{Finding $t$.} From $CE=EF$: $s^2+t^2=2(s-t)^2$, i.e. $t^2-4st+s^2=0$, so $t=(2-\sqrt3)s$ (the other root exceeds $s$). Hence $s-t=(\sqrt3-1)s$.

\textbf{Areas.} Triangle $DFG$ has base $DG=s-t$ and height $FG=s$, so $[DFG]=\dfrac{1}{2} s(s-t)=54$. Triangle $AEF$ is right-angled at $A$ with legs $s-t$, so
\[
[AEF]=\dfrac{1}{2}(s-t)^2=[DFG]\cdot\dfrac{s-t}{s}=54(\sqrt3-1).
\]
The area of triangle $AEF$ is $\boxed{54\sqrt3-54}\approx39.53$.
